% $Id: template.tex 11 2007-04-03 22:25:53Z jpeltier $
%
% CS 7389I -- Project Ideas
% TODO: Verify all citations on Google Scholar / ACM DL before final submission.
%       Each entry is noted with confidence level in template.bib.

\documentclass{styling/vgtc}                          % final (conference style)

\usepackage{mathptmx}
\usepackage{graphicx}
\usepackage{times}

\onlineid{0}
\vgtccategory{Research}
\vgtcinsertpkg

\title{Project Ideas: Novel Interaction Techniques in Immersive Environments}

\author{Andrew Scouten\\ %
        \scriptsize{yzb2@txstate.edu}\\ %
        \scriptsize Texas State University
\and Avery R Vanausdal\\ %
        \scriptsize{arv107@txstate.edu}\\ %
        \scriptsize Texas State University}

\abstract{This document presents five potential research project ideas for
CS~7389I, each exploring a novel interaction technique in virtual or mixed
reality. The ideas address AR-guided device repair, immersion-preserving
safety transitions, bystander awareness in VR, gaze-directed annotation in
mixed reality, and spatially-adaptive VR interface layout. Each idea is
motivated by a concrete research problem, defines a core interaction
contribution, proposes a study or prototype system to test it, and is
grounded in relevant prior literature.
} % end of abstract

\CCScatlist{
  \CCScat{H.5.2}{Information Interfaces and Presentation}{User Interfaces}%
{Interaction styles};
  \CCScat{H.5.1}{Information Interfaces and Presentation}{Multimedia Information Systems}%
{Artificial, augmented, and virtual realities}
}

%%%%%%%%%%%%%%%%%%%%%%%%%%%%%%%%%%%%%%%%%%%%%%%%%%%%%%%%%%%%%%%%
%%%%%%%%%%%%%%%%%%%%%% START OF THE PAPER %%%%%%%%%%%%%%%%%%%%%%
%%%%%%%%%%%%%%%%%%%%%%%%%%%%%%%%%%%%%%%%%%%%%%%%%%%%%%%%%%%%%%%%

\begin{document}

\firstsection{Promising Ideas}

\maketitle

From the five project ideas presented in this document, we identify the following three as
particularly noteworthy.

\subsection{Most Feasible Idea: DreamGuard}

DreamGuard is the most realistic to complete within the course timeline. The system relies
entirely on capabilities already exposed by the Meta Quest~3 SDK---color passthrough,
depth estimation, and person detection via the Passthrough Camera API---without requiring
custom hardware or difficult-to-recruit participant populations. The within-subjects
study design is straightforward to administer in a standard lab play space using a scripted
confederate, and the primary outcome measures (IPQ, SSQ, and motion-capture contact
logs) are well-established instruments that require no custom tooling to deploy.

\subsection{Most Creative Idea: SpatialFix}

SpatialFix represents the most unconventional concept in our set. Applying a learned
feature-matching pipeline (SuperGlue or LoFTR) to real-time passthrough video in order to
spatially anchor iFixIt repair instructions onto an arbitrary physical device---without
fiducial markers or CAD models---is a combination that has not been demonstrated on a
consumer headset. The idea reimagines the entire class of AR maintenance systems by
replacing per-device engineering effort with a generalized, appearance-based registration
approach, and pairs it with a hands-free, gaze- or voice-driven step-advancement interface
that keeps the technician's attention and hands fully on the repair task.

\subsection{Most Promising Idea: GazePin}

GazePin has the greatest potential to produce meaningful, generalizable research results.
The Midas Touch problem is well-known but remains unsolved for real-world object
annotation on consumer MR headsets, and the three-way empirical comparison (dwell,
gaze-plus-pinch, and pointing) maps directly onto a design question that headset
manufacturers and MR application developers face today. The study yields concrete,
actionable design guidelines supported by quantitative metrics (completion time, false
activation rate, NASA-TLX), and the availability of combined eye tracking and hand
tracking on the Meta Quest Pro makes it possible to run all three conditions on a single
commercially available device.

%%-------------------------------------------------------------------
\section{SpatialFix: AR-Guided Device Repair via Feature-Matched Instruction Overlays}

\subsection{Introduction}

Screen-based repair guides, such as those provided by iFixIt, require
technicians to constantly shift attention between an external display and the
physical device under repair, increasing cognitive load and the risk of
assembly errors. Spatial augmented reality affords the opportunity to anchor
step-by-step instructions directly onto the physical object, eliminating this
attention-splitting and freeing both hands for the task. A consumer passthrough
VR headset such as the Meta Quest~3---combining high-resolution color
passthrough, inside-out tracking, hand tracking, and an accessible developer
ecosystem---provides a compelling platform for exploring this interaction. Recent
reviews of AR-assisted maintenance also highlight the growing role of markerless
computer-vision pipelines in repair workflows~\cite{malta2023maintenance}.

\subsection{Core Contribution}

This project proposes a system that applies learned feature matching
(e.g., SuperGlue or LoFTR) to recognize a target device from the Quest
passthrough camera feed and spatially registers iFixIt repair steps as
anchored AR overlays onto the physical object in real time. The key
interaction contribution is a hands-free, gaze- or voice-driven
step-advancement interface that allows the technician to progress through
instructions without ever touching a secondary screen. Unlike prior AR
maintenance systems that required fiducial markers or CAD model registration,
the proposed approach relies entirely on appearance-based matching against
reference images scraped from publicly available repair guides, enabling the
technique to generalize to the breadth of the iFixIt device catalog without
per-device engineering effort.

\subsubsection{Reference Image Acquisition}
iFixIt publishes its repair guides under a Creative Commons license and provides a
public REST API that exposes guide metadata, step text, and step images for thousands
of devices. At project start we would query this API for a target set of devices
(e.g., a representative sample of laptop and smartphone models), downloading one or
more per-step reference images per device and organizing them into a local index keyed
by device identifier and step number. These reference images serve as the database
against which the feature-matching pipeline compares live passthrough frames at
runtime. No scraping of the iFixIt website itself would be necessary; the official API
endpoints support bulk retrieval within the terms of service. The resulting dataset would
be bundled with the prototype and made available alongside the project code repository.

\subsection{Proposed Study}

We would conduct a within-subjects comparative study in which
participants perform a standardized multi-step laptop disassembly and
reassembly task under three conditions: (1) a tablet displaying the iFixIt
web guide, (2) the proposed spatial AR overlay system on a Meta Quest~3, and
(3) an AR system using fiducial marker anchoring as a technical baseline.
Primary outcome measures are task completion time, number of assembly errors
(e.g., missed screws, incorrect component seating), and NASA-TLX workload
ratings. Secondary measures include head and hand motion patterns derived from
the headset sensors, as well as post-task semi-structured interview responses
probing trust in the spatial overlays and frustrations with feature-matching
failures or drift.

\subsection{Related Work}

Henderson and Feiner~\cite{henderson:2011:ar} conducted a seminal controlled
study comparing AR-overlaid maintenance documentation against paper and
monitor-based instructions for maintenance tasks inside an armored personnel carrier turret, finding
significant reductions in time-on-task and head movement when instructions
were spatially registered to the physical object. Our work extends this
finding to a self-contained consumer passthrough headset, replacing their
tethered see-through HMD and CAD-model registration pipeline with a
learned correspondence-based approach. Malta et al.~\cite{malta2023maintenance}
show that recent AR maintenance systems are increasingly moving toward
markerless computer-vision pipelines and more scalable deployment models,
which makes a photo-based registration strategy timely. By substituting
per-device CAD alignment with SuperGlue or LoFTR estimation, our approach
generalizes to arbitrary devices for which only a photographic repair guide
exists---a scalability advance that Henderson and Feiner's system did not
address.

%%-------------------------------------------------------------------
\section{DreamGuard: Immersion-Preserving Passthrough Safety Transitions in Virtual Reality}

\subsection{Introduction}

Consumer VR headsets such as the Meta Quest series rely on guardian grid
overlays to warn users of physical boundaries, but these grids abruptly
interrupt the virtual experience with a jarring visual signal that has no
aesthetic relationship to the virtual world. This break in immersion can
cause disorientation and, in extended-use scenarios, contribute to
cybersickness. Although current headsets increasingly support passthrough and
boundary visualization, recent work shows that open design questions remain
about how safety cues should communicate nearby hazards without unnecessarily
breaking immersion~\cite{qin2024brokenwall,wu2023guardian}.

\subsection{Core Contribution}

We propose \emph{DreamGuard}, a proximity-triggered safety system that
gradually blends the passthrough camera feed into the virtual scene using
aesthetically-matched visual transitions---soft fog, light bloom, or
environment-themed effects---rather than a hard grid overlay. The transition
is driven by depth estimation or person detection from the headset cameras,
begins before a collision is imminent, and fades back to full VR once the
user moves to safety. The core research contribution is the design and
empirical evaluation of immersion-preserving safety transitions: does a
seamless aesthetic blend produce fewer cybersickness symptoms and less
presence disruption than the standard guardian grid, while maintaining
equivalent safety outcomes?

\subsection{Proposed Study}

Participants would complete an immersive exploration task in a
standardized lab play space while scripted proximity events (a confederate
walking nearby, the user approaching a wall) occur at fixed intervals.
A within-subjects design would compare three conditions: (1) the standard
Meta Quest guardian grid, (2) DreamGuard with an aesthetic passthrough
blend, and (3) a no-warning control. Presence would be measured with the
Igroup Presence Questionnaire (IPQ) after each condition, cybersickness with
the Simulator Sickness Questionnaire (SSQ), and safety outcomes by logging
near-miss and contact events from motion-capture data.

\subsection{Related Work}

Milgram and Kishino~\cite{milgram:1994:taxonomy} introduced the
Reality-Virtuality Continuum, establishing that mixed-reality displays
exist on a spectrum between fully virtual and fully real environments and
that perceptual coherence across transitions along that spectrum is a
fundamental design challenge. DreamGuard operationalizes this continuum
dynamically: a safety event triggers a controlled traversal from fully
virtual toward the real end of the continuum, and recovery moves the user
back. Rather than a binary mode switch (as in the Guardian grid), our
system treats the transition itself as a designable interaction artifact.
Recent work on guardian cues and metaphor-based boundary transformations has
already shown that the form of a safety warning changes detection speed and
subjective experience~\cite{wu2023guardian,qin2024brokenwall}. Our project
extends that line of work by specifically testing gradual, aesthetically
integrated passthrough blends rather than discrete boundary widgets.

%%-------------------------------------------------------------------
\section{Who's There? Representing Nearby Bystanders in VR to Balance Collision Safety and Immersion}

\subsection{Introduction}

Standalone VR headsets still provide limited support for awareness of people
who walk into a user's play space mid-session, creating real risks of
collision and social awkwardness. Existing research and commercial systems
span full passthrough, abstract alerts, and blended awareness cues, but there
is still no clear design guidance on what person representation best
preserves immersion while reliably prompting collision avoidance~\cite{ohagan2023bystander,qin2024brokenwall}.

\subsection{Core Contribution}

This project investigates four progressively abstract representations of a
camera-detected bystander rendered inside the VR scene: (1)~a cropped
passthrough video silhouette, (2)~a stylized low-poly avatar, (3)~a ghosted
wireframe outline, and (4)~visible arms and hands only. The core contribution
is an empirical comparison of these conditions on both safety outcomes and
subjective immersion, producing concrete design guidelines for headset
manufacturers and VR developers. Unlike prior work focused on social presence
between co-present VR users, this work targets the asymmetric case of a
non-HMD bystander intruding into a single user's virtual space.

\subsection{Proposed Study}

Participants will complete an immersive VR task (a puzzle or navigation
game) in a lab space where a confederate bystander walks through the play
area at scripted intervals, with bystander representation assigned in a
within-subjects counterbalanced design across the four conditions. Primary
safety metrics will be near-miss and contact frequency derived from
motion-capture logs. Immersion and presence will be measured with the IPQ
after each condition, and social comfort and perceived intrusiveness will be
captured via post-hoc 7-point Likert scales. A secondary eye-tracking
measure will assess how often participants fixate on the bystander
representation, providing insight into the attentional cost of each style.

\subsection{Related Work}

Latoschik et al.~\cite{latoschik2017avatar} demonstrated that avatar
realism significantly affects social presence and behavior in immersive VR,
showing that even stylized representations can sustain strong social
signals. More recently, O'Hagan et al.~\cite{ohagan2023bystander} showed that
VR users' awareness needs during bystander interactions depend strongly on
social context rather than raw cue visibility alone. Together, these works
motivate our choice of a range of abstraction levels, from photorealistic
passthrough to minimalist hand outlines. Our project extends them by
addressing the asymmetric intrusion case of a non-HMD bystander whose
real-world geometry must be translated into a VR-native representation in
real time.

%%-------------------------------------------------------------------
\section{GazePin: Gaze-Directed Object Annotation in Mixed Reality with Multi-Modal Activation}

\subsection{Introduction}

Annotating or retrieving information about real-world objects in mixed
reality currently requires users to manually point or reach toward target
objects, which is physically effortful and fatiguing over time. Eye gaze is
a natural, low-effort pointing channel---users look at objects before
interacting with them---yet na\"{i}ve gaze-selection systems suffer from the
\emph{Midas Touch} problem: unintended activations occur because the system
cannot distinguish looking at an object from selecting it. No established
design solution balances the efficiency of gaze input with reliable,
intentional activation for real-world object annotation on consumer MR
headsets~\cite{pfeuffer2023palmgazer}.

\subsection{Core Contribution}

We propose \emph{GazePin}, a system that uses the Meta Quest Pro's eye
tracker\footnote{GazePin requires a Meta Quest Pro or other MR headset equipped with
an eye-tracking sensor. The standard Meta Quest~3 does not include eye tracking and
cannot run the gaze-based conditions; access to a Quest Pro (or equivalent device) is a
prerequisite for this project.} to detect fixation on a physical object via the passthrough camera
and triggers a contextual virtual annotation overlay using one of three
studied activation modalities: (1)~dwell time, (2)~gaze combined with a
hand pinch gesture, or (3)~hand pointing alone as a baseline. The core
research contribution is the first empirical comparison of these activation
methods for real-world object annotation on a consumer MR headset,
quantifying the tradeoff between annotation speed, false activation rate,
and cognitive load across modalities.

\subsection{Proposed Study}

A within-subjects study would ask participants to annotate a
series of labeled real-world objects in a structured lab environment under
the three activation conditions. Each condition requires selecting objects
of varying size and distance. Primary metrics are annotation completion
time, false activation rate (Midas Touch errors), and NASA-TLX cognitive
load scores. Secondary measures include subjective preference ratings and
semi-structured interview responses about perceived naturalness and
reliability of each activation method. Eye-tracking fixation data will be
logged throughout to analyze gaze patterns preceding correct and false
activations.

\subsection{Related Work}

Jacob~\cite{jacob:1991:eye} established the foundational paradigm for using
eye movements as a direct input channel in interactive systems and identified
the Midas Touch problem, proposing dwell time as a primary mitigation
strategy. More recent XR work such as PalmGazer~\cite{pfeuffer2023palmgazer}
suggests that eye--hand combinations may support selection in
AR interfaces, though its results were preliminary; it does not address annotation of physical scene objects.
Our project extends this line of work by moving from menu interaction to
consumer passthrough MR object annotation, introducing the challenge of
spatially registering gaze with physical real-world objects rather than
on-screen UI elements. The availability of combined eye tracking and hand
tracking on the Meta Quest Pro enables us to empirically evaluate the
gaze-plus-gesture condition with contemporary hardware.

%%-------------------------------------------------------------------
\section{RoomFit: Spatially-Adaptive VR Interface Layout Guided by Detected Room Geometry}

\subsection{Introduction}

Virtual reality interfaces---menus, panels, and toolbars---are typically
designed for a fixed virtual layout without accounting for the user's actual
physical room dimensions or the positions of walls and furniture. When
physical space is limited or irregularly shaped, users may reach toward a
virtual panel and collide with a real wall, or find themselves physically
constrained in ways that create frustration or break immersion. There is still limited design guidance for VR UIs that proactively adapt to
the physical environment the user inhabits, especially across rooms with
different sizes and furniture constraints~\cite{cheng2023interactionadapt}.

\subsection{Core Contribution}

We propose \emph{RoomFit}, a system that queries the Meta Quest's scene
understanding API to retrieve detected room geometry (walls, floor, ceiling,
and furniture bounding boxes) and uses this information to dynamically adapt
the placement, scale, and reach distance of VR UI panels at session start.
Panels are placed only in physically free space, interaction distances are
scaled to available arm reach within detected bounds, and UI elements that
would require the user to reach into a detected wall are relocated. The core
research contribution is an empirical evaluation of whether this spatial
adaptation reduces physical collisions and improves task performance
compared to a fixed-layout baseline, particularly in small or cluttered
rooms.

\subsection{Proposed Study}

A between-subjects study would have participants complete a series
of standardized VR menu navigation tasks in three physical room size
conditions (small: $2{\times}2$~m, medium: $3{\times}3$~m, large:
$4{\times}4$~m) under either the RoomFit adaptive layout or a fixed-layout
control. Primary metrics are the number of physical wall or furniture
contacts, task completion time, and presence as measured by the IPQ. We
hypothesize that the adaptive condition will show the greatest safety benefit
in small rooms, allowing us to characterize the conditions under which
spatial adaptation provides the most value.

The three room-size conditions would be realized within a single large lab
space rather than requiring three separate rooms. Colored tape on the floor
would mark the active play-area boundary for each condition, and portable
soft barriers (foam pool noodles or padded dividers on stands) would provide
physical feedback at the boundary edges. Before each condition block, the
Quest's Scene Understanding API would be re-invoked with the headset
positioned at the center of the newly marked zone so that the system ingests
the correct boundary geometry; furniture bounding boxes within the active
zone would be represented by cardboard box stand-ins of standardized sizes
placed at fixed relative positions. This approach keeps the physical setup
repeatable across participants while eliminating the logistical burden of
booking three different rooms.

\subsection{Related Work}

Lu et al.~\cite{lu2024adaptive} present the most directly related prior
work: an adaptive MR UI placement system that uses deep reinforcement
learning to position 3D interface panels relative to the user's detected
pose and physical environment, framing UI placement as an optimization
problem in the presence of dynamic real-world constraints. RoomFit shares
this core framing but targets a complementary scenario---room geometry known
at session start rather than pose-driven real-time repositioning---and
evaluates safety outcomes (physical contacts) rather than ergonomic
preferences. Cheng et al.~\cite{cheng2023interactionadapt} similarly showed
that VR workspace layouts can be adapted to physical affordances and
obstacles, confirming that environment-aware adaptation improves usability.
Together, these works establish that the physical context of a VR session is
a first-class design variable for interface layout; RoomFit extends this
direction to safety-oriented reconfiguration of standard menus and panels
driven by room-scan geometry.

%%-------------------------------------------------------------------

\acknowledgements{The authors would like to thank the CS~7389I course
instructors for their guidance on this project.}

\bibliographystyle{abbrv}
\bibliography{references/project_ideas}
\end{document}
